% Problem-section file following ../_template.tex.
\section{Unclassified existence parameters for homogeneous AME states}
\label{sec:problem-45}

\paragraph{Problem.}
For which of the parameter pairs specified below does an absolutely maximally
entangled state exist?  A normalized vector
$|\psi\rangle\in(\mathbb C^d)^{\otimes n}$ is an
$\operatorname{AME}(n,d)$ state when
\begin{equation}
  \operatorname{Tr}_{S^c}\!\left(|\psi\rangle\!\langle\psi|\right)
  =\frac{I_{d^{|S|}}}{d^{|S|}}
  \quad\text{for every }S\subseteq\{1,\ldots,n\}
  \text{ with }|S|\leq\left\lfloor\frac n2\right\rfloor.
  \label{eq:p45-ame-condition}
\end{equation}
Determine whether a state satisfying Eq.~\eqref{eq:p45-ame-condition} exists
for every pair in
\begin{equation}
  \mathcal U:=\bigl\{(8,6),(8,10),(9,6),(9,10),(10,6),(10,10),
  (11,3),(11,6),(11,10),(12,6),(12,10)\bigr\}.
  \label{eq:p45-unclassified-pairs}
\end{equation}
Thus the task is to classify every pair in Eq.~\eqref{eq:p45-unclassified-pairs}
by existence or nonexistence, without restricting to stabilizer or
minimal-support states.

\subsection*{Status}

Unsolved

\subsection*{Source}

The unresolved parameter list is drawn from the AME existence table and open
tasks compiled by Rajchel-Mieldzio\'c et al., with cases removed when later
constructions settled them \sourcecite{ref:p45-review}{RBR+26}.

\subsection*{Progress}

\begin{itemize}
  \item Existence under Eq.~\eqref{eq:p45-ame-condition} is equivalent to a
  pure one-dimensional quantum error-correcting code with parameters
  $((n,1,\lfloor n/2\rfloor+1))_d$.  Quantum-code bounds and shadow
  inequalities therefore provide nonexistence tests, but they do not settle
  the pairs in Eq.~\eqref{eq:p45-unclassified-pairs}
  \sourcecite{ref:p45-scott}{Sco04},
  \sourcecite{ref:p45-huber-grassl}{HG20}.

  \item General constructions yield $\operatorname{AME}(5,d)$ and
  $\operatorname{AME}(6,d)$ for every $d\geq2$, while the qubit existence
  problem is completely classified: an $\operatorname{AME}(n,2)$ state
  exists exactly for $n\in\{2,3,5,6\}$.  These results delimit, but do not
  decide, Eq.~\eqref{eq:p45-unclassified-pairs}
  \sourcecite{ref:p45-review}{RBR+26}.

  \item Two 2026 advances remove stale benchmark cases: the seven-party
  classification proves $\operatorname{AME}(7,d)$ exists exactly when
  $d\geq3$, and an exact Hermitian self-dual MDS code constructs
  $\operatorname{AME}(12,5)$.  Neither result decides any pair retained in
  Eq.~\eqref{eq:p45-unclassified-pairs}
  \sourcecite{ref:p45-shi}{SZZL26},
  \sourcecite{ref:p45-bevins-bidav}{BB26}.
\end{itemize}

\subsection*{References}

\begin{enumerate}[label={},leftmargin=4.8em,labelsep=0.5em]
  \footnotesize
  \item[\textup{[Sco04]}]\label{ref:p45-scott}
  A. J. Scott,
  ``Multipartite Entanglement, Quantum-Error-Correcting Codes, and Entangling
  Power of Quantum Evolutions,'' \emph{Physical Review A} \textbf{69},
  052330 (2004).
  \href{https://doi.org/10.1103/PhysRevA.69.052330}%
       {doi:10.1103/PhysRevA.69.052330};
  \href{https://arxiv.org/abs/quant-ph/0310137}{arXiv:quant-ph/0310137}.

  \item[\textup{[HG20]}]\label{ref:p45-huber-grassl}
  F. Huber and M. Grassl,
  ``Quantum Codes of Maximal Distance and Highly Entangled Subspaces,''
  \emph{Quantum} \textbf{4}, 284 (2020).
  \href{https://doi.org/10.22331/q-2020-06-18-284}%
       {doi:10.22331/q-2020-06-18-284};
  \href{https://arxiv.org/abs/1907.07733}{arXiv:1907.07733}.

  \item[\textup{[RBR+26]}]\label{ref:p45-review}
  G. Rajchel-Mieldzio{\'c}, R. Bistro{\'n}, A. Rico, A. Lakshminarayan, and
  K. {\.{Z}}yczkowski,
  ``Absolutely Maximally Entangled Pure States of Multipartite Quantum
  Systems,'' \emph{Reports on Progress in Physics} \textbf{89}, 057601
  (2026).
  \href{https://doi.org/10.1088/1361-6633/ae672a}%
       {doi:10.1088/1361-6633/ae672a};
  \href{https://arxiv.org/abs/2508.04777}{arXiv:2508.04777}.

  \item[\textup{[SZZL26]}]\label{ref:p45-shi}
  F. Shi, X. Zhang, Q. Zhao, and L. Li,
  ``Complete Existence Classification of Seven-Partite Absolutely Maximally
  Entangled States,'' arXiv:2608.01011 (2026).
  \href{https://doi.org/10.48550/arXiv.2608.01011}%
       {doi:10.48550/arXiv.2608.01011};
  \href{https://arxiv.org/abs/2608.01011}{arXiv:2608.01011}.

  \item[\textup{[BB26]}]\label{ref:p45-bevins-bidav}
  S. Bevins and Y. Bidav,
  ``Symmetry-Guided Constructions of Absolutely Maximally Entangled States in
  Five Open Cases,'' arXiv:2608.05781v2 (2026).
  \href{https://doi.org/10.48550/arXiv.2608.05781}%
       {doi:10.48550/arXiv.2608.05781};
  \href{https://arxiv.org/abs/2608.05781}{arXiv:2608.05781}.
\end{enumerate}

\subsection*{Comment}

The set in Eq.~\eqref{eq:p45-unclassified-pairs} is the nonredundant remainder
of the audited DeepMind benchmark list: the solved pairs $(7,6)$, $(7,10)$,
and $(12,5)$ are excluded, while the unresolved $(8,4)$ case is already
Problem 40.  Problems 41--43 respectively impose a real-coefficient
constraint, ask for minimal support, and classify local-unitary orbits; none
duplicates the unrestricted existence predicates collected here.

\subsection*{Tag}

Absolutely maximally entangled states; Quantum error correction; Multipartite
entanglement; Qudit systems; Quantum coding theory; Combinatorics

\subsection*{ID}

\texttt{op\_439ae5e7e9b3b043}
